\documentclass[11pt]{article}
\usepackage[margin=1in]{geometry}
\usepackage{amsmath,amssymb,amsthm}
\title{Disproof of Conjecture 00000004761}
\author{earthking11}
\date{October 6, 2026}
\newtheorem{proposition}{Proposition}
\begin{document}
\maketitle
\section*{Counterexample}
Let $NC(1)$ be the lattice of noncrossing partitions of a one-element set.
It has the single element $x=\{\{1\}\}$.  The defining convolution-inverse
identity for the Mobius function forces
\[
 \mu(x,x)=1.
\]
Indeed, on a one-element interval the zeta function is $\zeta(x,x)=1$, so its
convolution inverse must also take the value $1$.

The total Mobius sum over this lattice is consequently
\[
 \sum_{u\in NC(1)}\mu(x,u)=\mu(x,x)=1\ne0.
\]

\begin{proposition}
The total Mobius sum on a noncrossing-partition lattice is not always zero.
\end{proposition}
\begin{proof}
$NC(1)$ gives the one-term sum $1$ displayed above.
\end{proof}

The standard zero-sum formula is
\[
 \sum_{x\le z\le y}\mu(x,z)=0 \quad\text{when }x<y.
\]
For the degenerate interval $x=y$, the same sum equals $1$.  Thus a strictness
hypothesis repairs this part of the statement, but that hypothesis is absent
from the filed conjecture.  References to parity or nested cycles cannot alter
the diagonal normalization required by Mobius inversion.

The accompanying Lean project formalizes the one-element lattice, verifies
the inverse equation, and proves that the asserted universal zero total fails.
\end{document}
